%%%%%%%%%%%%%%%%%%%%%%%%%%%%%%%%%%%%%%%%%%%%%%%%%%%%%%%%%%%%%%%%%%%%%%
% Lớp 11 | Xác suất | Biến cố hợp và quy tắc cộng xác suất
% Dạng 1D9?2-1: Câu hỏi lý thuyết
% Số câu: 13  (sinh tự động bởi tach_cau.py; từ nay thầy sửa trực tiếp file này)
%%%%%%%%%%%%%%%%%%%%%%%%%%%%%%%%%%%%%%%%%%%%%%%%%%%%%%%%%%%%%%%%%%%%%%

%%%%%%%%%%%%%%%% TRẮC NGHIỆM 4 PHƯƠNG ÁN %%%%%%%%%%%%%%%%

% nguồn: Ok_TracNghiemLop11_Lan1
\begin{ex}%[1D9N2-1]%[ma:13e148ab]
	Cho hai biến cố A và B. Biến cố hợp của A và B là  biến cố
	\choice 
		{\lq\lq  A và B xảy ra\rq\rq }
		{\True \lq\lq  A hoặc B xảy ra\rq\rq }
		{\lq\lq A  xảy ra\rq\rq }
		 {\lq\lq B xảy ra hoặc cả A và B xảy ra\rq\rq }
	 \loigiai{
	Dựa  vào định nghĩa.}
\end{ex}

% nguồn: Ok_TracNghiemLop11_Lan1
\begin{ex}%[1D9N2-1]%[ma:edab2ef9]
	Cho $A$, $B$ là hai biến cố  xung khắc. Đẳng thức nào sau đây đúng?
	\choice 
		{\True $\mathrm{P}(A\cup  B)=\mathrm{P}(A)+\mathrm{P}(B)$}
		{$\mathrm{P}(A\cup B)=\mathrm{P}(A)\cdot \mathrm{P}(B)$} 
		{$\mathrm{P}(A\cup  B)=\mathrm{P}(A)-\mathrm{P}(B)$}
		{$\mathrm{P}(A\cap B)=\mathrm{P}(A)+\mathrm{P}(B)$} 
	\loigiai{ 
		Ta có $\mathrm{P}(A\cup  B)=\mathrm{P}(A)+\mathrm{P}(B)-\mathrm{P}(A\cap B)$.\\
		Vì  $A$, $B$ là hai biến cố xung khắc nên $A\cap B=\varnothing $.\\
	Từ đó  suy ra $\mathrm{P}(A\cup B)=\mathrm{P}(A)+\mathrm{P}(B)$.}
\end{ex}

% nguồn: Ok_TracNghiemLop11_Lan1
\begin{ex}%[1D9N2-1]%[ma:28278385]
	Hai bạn Triết và Khang, mỗi người gieo đồng thời một con xúc xắc cân đối, đồng chất. Xét hai  biến cố sau:\\
	A: \lq\lq Số chấm xuất hiện trên con xúc xắc bạn Triết gieo là số  lẻ.\rq\rq .\\
	B: Số chấm xuất hiện trên con  xúc xắc bạn Khang gieo là số chia hết cho 5.\rq\rq .\\
	 Chọn câu đúng:
	\choice 
		{\True Hai biến cố A, B là hai biến  cố độc lập}
		{Hai  biến cố A, B là hai biến cố xung khắc}
		 {Biến cố A giao B là $\{1;3;5\}$}
		{Biến cố A hợp B là $\{1;2;3;5\}$} 
	 \loigiai{
	Dựa vào  định nghĩa.}
\end{ex}

% nguồn: Ok_TracNghiemLop11_Lan1
\begin{ex}%[1D9N2-1]%[ma:84235060]
	Cho hai biến cố $A$ và $B$ có $\mathrm{P}(A)=\dfrac{1}{3}$,  $\mathrm{P}(B)=\dfrac{1}{4}$, $\mathrm{P}(A\cup B)=\dfrac{1}{2}$. Khẳng định nào sau đây đúng
	 \choice
		{Độc  lập}
		{\True Không xung khắc} 
		 {Xung khắc}
		{Không rõ} 
	 \loigiai{
		Ta xét 
		\begin{eqnarray*}
			&  &\mathrm{P}(A\cup B)= \mathrm{P}(A)+P(B)-P(A\cap B)\\
			&\Leftrightarrow & \dfrac{1}{2}=\dfrac{1}{3}+\dfrac{1}{4}-\mathrm{P}(A\cap B)\\ 
			&\Leftrightarrow &  \mathrm{P}(A\cap B)=\dfrac{1}{12}\ne 0.
		\end{eqnarray*}
	Suy  ra hai biến cố $A$ và $B$ là hai biến cố không xung khắc.}
\end{ex}

% nguồn: Ok_TracNghiemLop11_Lan1
\begin{ex}%[1D9N2-1]%[ma:071a4d1d]
	Gieo $2$ con xúc xắc cân đối và đồng chất. Gọi  $A$ là biến cố \lq\lq Tích số chấm xuất hiện là số lẻ\rq\rq. Biến cố nào sau đây xung khắc với biến cố $A$?
	\choice 
		{\lq\lq Xuất hiện hai  mặt có cùng số chấm\rq\rq}
		{\True  \lq\lq Tổng số chấm xuất hiện là số lẻ\rq\rq}
		{\lq\lq Xuất hiện it nhất một mặt  có số chấm là số lẻ\rq\rq}
		{\lq\lq Xuất hiện hai  mặt có số chấm khác nhau\rq\rq}
	 \loigiai{
	Ta có biến cố \lq\lq Tổng số chấm xuất hiện  là số lẻ\rq\rq\,xung khắc với biến cố \lq\lq Tích số chấm xuất hiện là số lẻ\rq\rq.}
\end{ex}

% nguồn: Ok_TracNghiemLop11_Lan1
\begin{ex}%[1D9N2-1]%[ma:bcc6dd80]
	Hai bạn  Dương và Khoa, mỗi người gieo đồng thời một con xúc xắc cân đối, đồng chất. Xét hai biến cố sau:\\
	E: \lq\lq Số chấm xuất hiện trên con  xúc xắc bạn Dương gieo là số chính phương\rq\rq .\\
	B: \lq\lq Số chấm xuất hiện trên con xúc xắc của bạn  Khoa là số chia hết cho 4\rq\rq .\\
	Chọn  câu đúng:
	\choice 
		{\True Hai biên cố E  và B độc lập}
		{Hai biến cố E và B không độc  lập}
		{Hai biến cố E và  B xung khắc}
		{Hai biến cố E và B không  xung khắc}
	 \loigiai{
	 Dựa vào định nghĩa}
\end{ex}

% nguồn: Ok_TracNghiemLop11_Lan1
\begin{ex}%[1D9N2-1]%[ma:cbe57852]
	Một hộp có $12$ chiếc thẻ cùng loại, mỗi thẻ được ghi một trong các số tự nhiên,  hai thẻ khác nhau thì ghi hai số khác nhau. Rút ngẫu nhiên 1 chiếc thẻ trong hộp. Xét biến cố A: \rq\rq Số xuất hiện trên thẻ được rút ra là số chia hết cho 3\lq\lq\, và biến cố B: \lq\lq Số xuất hiện trên thẻ được rút ra là số chia hết cho 4\rq\rq.\\
	Biến cố A  hợp B được phát biểu như sau:
	 \choice
		{\lq\lq Số xuất hiện trên thẻ là số vừa chia hết cho 3 vừa chia hết cho 4\rq\rq  }
		{\True \lq\lq Số xuất hiện trên thẻ là số chia hết cho 3 hoặc chia  hết cho 4\rq\rq }
		{\lq\lq Số xuất hiện  trên thẻ là số chia hết cho 12\rq\rq }
		{\lq\lq Số xuất hiện trên thẻ  là số chia hết cho 7\rq\rq }
	\loigiai{ 
	Dựa  vào định nghĩa.}
\end{ex}

% nguồn: Ok_TracNghiemLop11_Lan1
\begin{ex}%[1D9H2-1]%[ma:0008c101]
	Nếu $A$ và $B$  là hai biến cố thì $\mathrm{P}(A \cup B)$ bằng
	\choice 
		{\True  $\mathrm{P}(A)+\mathrm{P}(B)-\mathrm{P}(A \cap B)$}
		 {$\mathrm{P}(A)-\mathrm{P}(B)-\mathrm{P}(A \cap B)$}
		 {$\mathrm{P}(A), \mathrm{P}(B)-\mathrm{P}(A \cap B)$}
		{$\mathrm{P}(A)  \cdot \mathrm{P}(B)+\mathrm{P}(A \cap B)$}
	 \loigiai{Vì $A$ và $B$ là hai biến cố bất kì nên $\mathrm{P}(A\cup B)=\mathrm{P}(A)+\mathrm{P}(B)-\mathrm{P}(A \cap B)$.
	 }
\end{ex}


%%%%%%%%%%%%%%%% ĐÚNG/SAI %%%%%%%%%%%%%%%%

% nguồn: Ok_dungsai_lop11_Lan1
\begin{ex}%[1D9N2-1]%[ma:331b9264]
 Cho hai biến cố $A$ và $B$. Hãy cho biết tính đúng sai của các mệnh  đề sau.
  \choiceTF
  {\True Nếu $\mathrm{P}(A  B) \neq \mathrm{P}(A) \mathrm{P}(B)$ thì hai biến cố $A$ và $B$ không độc lập}
  {\True Nếu hai biến cố $A$ và $B$ độc lập thì  $\mathrm{P}(A B)=\mathrm{P}(A) \mathrm{P}(B)$}
  {\True Cho hai biến cố xung khắc $A$ và  $B$. Khi đó $\mathrm{P}(A \cup B)=\mathrm{P}(A)+\mathrm{P}(B)$}
  {\True  Cho hai biến cố $A$ và $B$. Khi đó $\mathrm{P}(A \cup B)=\mathrm{P}(A)+\mathrm{P}(B)-\mathrm{P}(A B)$}
  \loigiai{
  Cho hai biến cố $A$  và $B$.
  \begin{itemchoice}
   \itemch Nếu $\mathrm{P}(A B) \neq \mathrm{P}(A) \mathrm{P}(B)$ thì hai biến cố $A$ và $B$  không độc lập.
   \itemch Nếu hai biến cố $A$ và $B$ độc lập thì  $\mathrm{P}(A B)=\mathrm{P}(A) \mathrm{P}(B)$.
   \itemch  Cho hai biến cố xung khắc $A$ và $B$. Khi đó $\mathrm{P}(A \cup B)=\mathrm{P}(A)+\mathrm{P}(B)$.
   \itemch Cho hai biến cố $A$ và $B$. Khi đó $\mathrm{P}(A \cup  B)=\mathrm{P}(A)+\mathrm{P}(B)-\mathrm{P}(A B)$.
  \end{itemchoice}
 } 
\end{ex}

% nguồn: Ok_dungsai_lop11_Lan1
\begin{ex}%[1D9N2-1]%[ma:cce81199]
 Cho hai biến cố $A$ và $B$. 
 \choiceTF 
  {\True Biến cố \lq\lq$A$ hoặc $B$ xảy ra\rq\rq, kí hiệu là $A\cup B$ được gọi  là \textbf{\textit{biến cố hợp}} của $A$ và $B$}
  {\True Biến cố $A \cup B$ xảy ra khi có ít nhất một trong hai biến cố  $A$ và $B$ xảy ra}
  {Tập hợp mô tả biến cố $A \cap  B$ là hợp của hai tập hợp mô tả biến cố $A$ và biến cố $B$}
  {\True Biến cố $A$ và biến cố $B$ được gọi là xung  khắc nếu $A$ và $B$ không đồng thời xảy ra}
 \loigiai{ 
   Cho hai biến cố $A$ và $B$.
  \begin{itemchoice}
   \itemch Biến cố \lq\lq$A$ hoặc  $B$ xảy ra\rq\rq, kí hiệu là $A\cup B$ được gọi là \textbf{\textit{biến cố hợp}} của $A$ và $B$.
   \itemch Biến cố $A \cup B$ xảy ra khi  có ít nhất một trong hai biến cố $A$ và $B$ xảy ra.
   \itemch Tập hợp mô tả biến cố $A \cup B$ là hợp của hai tập hợp mô tả biến cố $A$ và biến  cố $B$.
   \itemch Biến cố $A$ và biến cố $B$ được gọi  là xung khắc nếu $A$ và $B$ không đồng thời xảy ra.
  \end{itemchoice}
  }
\end{ex}

% nguồn: Ok_dungsai_lop11_Lan1
\begin{ex}%[1D9N2-1]%[ma:3c8ad6d1]
 Gieo một con xúc xắc cân đối, đồng chất liên tiếp hai lần. Xét các  biến cố sau \\
 \indent \indent $A\colon$ \lq\lq Số  chấm xuất hiện trong hai lần gieo đều là số chẵn\rq\rq;\\
 \indent \indent $B\colon$ \lq\lq Số chấm xuất hiện trong hai lần gieo khác tính chẵn  lẻ\rq\rq;\\
 \indent \indent $C\colon$ \lq\lq Tích số chấm xuất hiện trong hai lần gieo là số  chẵn\rq\rq;\\
 \indent \indent $D\colon$ \lq\lq Tổng số chấm xuất hiện trong hai lần gieo là số lớn  hơn $9$\rq\rq.\\
 Các khẳng định sau  đúng hay sai?
 \choiceTF 
  {\True Xác suất của biến  cố $A$ là $\dfrac{1}{4}$}
  {\True Biến cố $C$ là hợp của hai biến  cố $A$ và $B$}
  {Xác suất  của biến cố $C$ là $\dfrac{2}{3}$}
  {Xác suất của biến  cố $D$ là $\dfrac{3}{4}$}
 \loigiai{ 
  \begin{itemchoice}
   \itemch Ta có $n\left(\Omega \right)=6\cdot  6=36$; $n\left(A \right)=3\cdot 3=9$.\\
    Suy ra $\mathrm{P}\left(A \right)=\dfrac{n\left(A \right)}{n\left(\Omega \right)}=\dfrac{9}{36}=\dfrac{1}{4}$.
   \itemch Biến cố $C$ xảy ra khi và chỉ khi trong hai lần gieo có ít nhất một lần gieo  xuất hiện mặt có số chấm là số chẵn.\\
   Nếu cả hai lần gieo đều ra số chẵn thì biến cố  $A$ xảy ra.\\
   Nếu chỉ có một lần gieo ra số chẵn thì  biến cố $B$ xảy ra.\\
   Vậy biến cố $C$  là hợp của hai biến cố $A$ và $B$.
   \itemch Ta có $\overline{C}\colon$ \lq\lq Tích số chấm xuất hiện trong hai lần gieo là số lẻ\rq\rq; tức là \lq\lq Số chấm xuất hiện trong hai lần gieo đều  là số lẻ\rq\rq. Suy ra $n\left(\overline{C} \right)=3\cdot 3=9$.\\
   Do đó $\mathrm{P}\left(\overline{C} \right)=\dfrac{n\left(\overline{C} \right)}{n\left(\Omega \right)}=\dfrac{9}{36}=\dfrac{1}{4}\Rightarrow \mathrm{P}\left(C \right)=1-\mathrm{P}\left(\overline{C} \right)=1-\dfrac{1}{4}=\dfrac{3}{4}$. 
    \itemch Tổng số chấm xuất hiện trong hai lần gieo là số lớn hơn $9$ chỉ có các trường hợp sau
   $$\left(4;6\right);\left(5;5\right);\left(5;6\right);\left(6;4\right);\left(6;5\right);\left(6;6\right).$$ 
   Do đó $n\left(D \right)=6$. Suy ra $\mathrm{P}\left(D \right)=\dfrac{n\left(D  \right)}{n\left(\Omega \right)}=\dfrac{6}{36}=\dfrac{1}{6}$.
  \end{itemchoice}
  }
\end{ex}

% nguồn: Ok_dungsai_lop11_Lan1
\begin{ex}%[1D9H2-1]%[ma:9fe22b0b]
  Một hộp có $20$ chiếc thẻ cùng loại, được đánh số từ $1,2,\ldots,19,20$. Rút ngẫu nhiên một thẻ trong hộp. Xét các biến cố:\\
 $A:$ \lq\lq Số trên thẻ được rút ra là số chia  hết cho $5$\rq\rq.\\
 $B:$ \lq\lq Số trên thẻ  được rút ra là số chia hết cho $3$. \rq\rq\\
 $C:$ \lq\lq Số trên thẻ được rút ra là số chia hết cho $5$ hoặc chia hết cho $3$. \rq\rq\\ 
 $D:$ \lq\lq Số trên  thẻ được rút ra là số chia hết cho $15$ \rq\rq. \\
 Khi  đó:
  \choiceTF
  {Biến cố $A$ và  biến cố $D$ là hai biến cố xung khắc}
  {\True  Xác suất của biến cố $D$ là $\dfrac{1}{20}$}
  {Biến cố $D$ là biến cố hợp của biến cố $A$  và biến cố $B$}
  {Số phần tử của không gian mẫu là $190$} 
 \loigiai{ 
  \begin{itemize}
   \item Vì biến cố $D$ xảy ra thì biến cố $A$ cũng xảy ra nên hai biến cố  $A$ và $D$ không xung khắc.
   \item  Ta có $\mathrm{P}(D)=\dfrac{1}{20}$.
   \item Biến cố $D$ là biến cố giao của biến  cố $A$ và biến cố $B$.
   \item  Số phần tử không gian mẫu $n(\Omega)=20$.
  \end{itemize}
 } 
\end{ex}

% nguồn: Ok_dungsai_lop11_Lan1
\begin{ex}%[1D9H2-1]%[ma:e63ebedb]
 Một hộp có $20$ chiếc thẻ cùng loại, mỗi thẻ được ghi một  số trong các số {$1,2,\ldots,19,20$}; hai thẻ khác nhau thì ghi hai số khác nhau. Rút ngẫu nhiên một chiếc thẻ trong hộp. Xét các biến cố\\
 \indent $A\colon$ \lq\lq Số  trên thẻ được rút ra là số chia hết cho $5$\rq\rq.\\
 \indent $B\colon$ \lq\lq Số trên thẻ được rút ra là số chia hết  cho $3$\rq\rq.\\
 \indent $C\colon$ \lq\lq Số trên thẻ được rút ra là số chia hết  cho $5$ hoặc chia hết cho $3$\rq\rq.\\
 \indent $D\colon$ \lq\lq Số trên thẻ được rút ra là  số chia hết cho $15$\rq\rq.
  \choiceTF
  {Biến cố $A$ và biến cố $D$ là hai  biến cố xung khắc}
   {\True Xác suất của biến cố $D$ là $\dfrac{1}{20}$}
  {Biến cố $D$ là biến cố hợp của biến cố $A$ và biến  cố $B$}
  {Số phần  tử của không gian mẫu là $190$}
  \loigiai{
  \begin{itemchoice}
   \itemch Ta có thẻ số $15$ vừa chia hết cho cả $3$ và  $5$ nên $15$ vừa là phần tử của biến cố $A$, vừa là phần tử của biến cố $B$ nên hai biến cố $A$, $B$ không xung khắc.
   \itemch Xác suất của biến cố $D$  là $\mathrm{P}(D)=\dfrac{1}{20}$.
   \itemch Ta  có $D=A\cap B$.
   \itemch Số phần tử của  không gian mẫu là $n(\Omega)=\mathrm{C}_{20}^1=20$.
  \end{itemchoice}
 } 
\end{ex}
